\documentclass[a4paper]{article}
\usepackage{modulITERA}
\begin{document}
\SampulDok{RPS}{Rencana Pembelajaran Semester}{Capaian, rencana 16 pertemuan, penilaian, dan rubrik. Disusun dari sumber data yang sama dengan slide, modul, dan hands-on.}
\par\vspace{4pt}\noindent{\fontsize{9.5pt}{13pt}\selectfont
\begin{tabular}{@{}>{\raggedright\arraybackslash}p{0.280\textwidth}>{\raggedright\arraybackslash}p{0.700\textwidth}@{}}
\kepala{Komponen} & \kepala{Keterangan}\\
Kode dan nama & IF25-11002 Praktikum Pemrograman\\\garis
Bobot & 2 SKS praktikum (16 pertemuan × 150 menit)\\\garis
Semester & 1 (Ganjil)\\\garis
Prasyarat & tidak ada\\\garis
Mata kuliah pasangan & IF25-11001 Algoritma Pemrograman, topik sejajar setiap minggu\\\garis
Program studi & S1 Teknik Informatika, Fakultas Teknologi Industri, Institut Teknologi Sumatera\\\garis
Tim pengajar & Tim Pengajar Praktikum Pemrograman\\\garis
Bahasa dan alat & C++ (standar C++17); GDB Online di kelas, g++ di komputer sendiri\\\garis
Tahun & 2026\\
\garisdasar
\end{tabular}}\par\vspace{6pt}
\newpage
\Bagian{RPS}{Deskripsi mata kuliah}
Praktikum Pemrograman melatih mahasiswa menerjemahkan algoritma menjadi program C++ yang berjalan, lalu menjelaskan mengapa program itu berjalan. Materi mencakup struktur program, tipe data dan masukan, operator, percabangan, perulangan, array satu dan dua dimensi, string, pencarian, pengurutan, fungsi, rekursi, dan struct. Setiap pertemuan reguler memakai siklus problem-first lima fase: review, struggle, materi, hands-on, dan Exit Ticket.

\Bagian{RPS}{Capaian pembelajaran}
\par\vspace{4pt}\noindent{\fontsize{9.5pt}{13pt}\selectfont
\begin{tabular}{@{}>{\raggedright\arraybackslash}p{0.160\textwidth}>{\raggedright\arraybackslash}p{0.820\textwidth}@{}}
\kepala{Kode} & \kepala{Rumusan}\\
CPL06 & Mampu menganalisis dan mengimplementasi konsep dasar matematika, statistika, algoritma dan sistem komputer sebagai fondasi komputasional.\\\garis
CPMK0607 & Mahasiswa mampu menerapkan konsep dasar algoritma dan sistem komputer untuk menyelesaikan masalah komputasi. Diukur lewat Bagian A setiap instrumen.\\\garis
CPMK0608 & Mahasiswa mampu menjelaskan konsep dasar algoritma dan sistem komputer untuk menyelesaikan masalah komputasi. Diukur lewat Bagian B setiap instrumen.\\
\garisdasar
\end{tabular}}\par\vspace{6pt}
\Sub{Sub-CPMK per pertemuan}
Setiap pertemuan menurunkan CPMK0607 menjadi Sub-CPMK n.1 (menerapkan) dan CPMK0608 menjadi Sub-CPMK n.2 (menjelaskan).
\par\vspace{4pt}\noindent{\fontsize{8.5pt}{11.5pt}\selectfont
\begin{longtable}{@{}>{\raggedright\arraybackslash}p{0.060\textwidth}>{\raggedright\arraybackslash}p{0.460\textwidth}>{\raggedright\arraybackslash}p{0.460\textwidth}@{}}
\kepala{P} & \kepala{Sub-CPMK n.1 (CPMK0607)} & \kepala{Sub-CPMK n.2 (CPMK0608)}\\
\endhead
1 & Menulis, mengompilasi, dan menjalankan program C++ yang menampilkan keluaran terformat dengan \texttt{cout}, \texttt{endl}, dan escape sequence & Menjelaskan peran setiap bagian program C++ dan alur dari kode sumber sampai program berjalan, termasuk membaca pesan error kompilasi\\\garis
2 & Menulis program yang menyimpan data dalam variabel bertipe tepat (\texttt{int}, \texttt{double}, \texttt{char}, \texttt{string}, \texttt{bool}), membaca masukan dengan \texttt{cin}, dan menampilkan hasil ekspresi & Menjelaskan pemilihan tipe data dan menelusuri nilai variabel setelah serangkaian penugasan, termasuk akibat pembagian bilangan bulat\\\garis
3 & Menulis program yang memakai operator aritmetika termasuk \texttt{\%}, operator penugasan majemuk, increment, serta operator relasional dan logika untuk menyelesaikan perhitungan & Menjelaskan urutan evaluasi sebuah ekspresi berdasarkan precedence dan asosiativitas, serta nilai yang dihasilkan pada setiap langkah\\\garis
4 & Menulis program yang memilih tindakan dengan \texttt{if}, \texttt{if-else}, rantai \texttt{else if}, \texttt{switch}, dan operator ternary sesuai syarat masalah & Menjelaskan cabang mana yang dijalankan untuk masukan tertentu dan merancang data uji yang mencakup setiap cabang dan nilai batasnya\\\garis
5 & Menulis program yang mengulang proses dengan \texttt{for}, \texttt{while}, dan \texttt{do-while}, memakai pola pencacah, akumulator, dan nilai sentinel & Menelusuri perulangan dengan trace table, menentukan berapa kali badan perulangan dijalankan, dan menjelaskan penyebab off-by-one atau perulangan tak berhenti\\\garis
6 & Menulis program dengan perulangan bersarang, \texttt{break}, dan \texttt{continue} untuk menghasilkan pola dan memproses data dua dimensi & Menurunkan hubungan antara nomor baris dan banyak isi pada sebuah pola, serta menelusuri perulangan bersarang untuk menentukan keluarannya\\\garis
7 & Menyelesaikan masalah yang menggabungkan masukan, ekspresi, percabangan, dan perulangan dalam satu program C++ dalam batas waktu ujian & Menelusuri program yang menggabungkan percabangan dan perulangan bersarang, lalu menjelaskan setiap perubahan nilai dengan istilah yang tepat\\\garis
8 & Menyelesaikan masalah komputasi dengan program C++ yang memakai variabel, ekspresi, percabangan, dan perulangan & Menelusuri dan menjelaskan perilaku program yang memakai percabangan dan perulangan\\\garis
9 & Menulis program yang menyimpan data sejenis di array satu dimensi dan memprosesnya dengan perulangan: mengisi, menampilkan, menjumlahkan, mencari maksimum, dan menggeser & Menelusuri isi array sesudah setiap langkah pemrosesan dan menjelaskan kesalahan indeks, termasuk akses di luar batas\\\garis
10 & Menulis program yang memproses data tabel dengan array dua dimensi dan mengolah teks dengan \texttt{string}: mengakses karakter, menghitung, mengubah huruf, dan memakai fungsi anggota \texttt{string} & Menelusuri indeks baris dan kolom yang diakses pada perulangan bersarang, serta menjelaskan pemrosesan string karakter demi karakter\\\garis
11 & Menulis program yang mencari data pada array dengan linear search dan binary search, termasuk mencari semua kemunculan dan menangani data yang tidak ditemukan & Menelusuri langkah binary search dengan trace table low, high, dan mid, serta menjelaskan perbedaan banyak perbandingan kedua algoritma\\\garis
12 & Menulis program yang mengurutkan array naik maupun turun dengan bubble sort dan selection sort, termasuk mengurutkan dua array sejajar sekaligus & Menelusuri isi array sesudah setiap putaran pengurutan dan menjelaskan banyak perbandingan serta pertukaran yang terjadi\\\garis
13 & Menulis program yang dipecah menjadi fungsi dengan parameter dan nilai kembali yang tepat, termasuk fungsi \texttt{void}, fungsi \texttt{bool}, parameter referensi, dan parameter array & Menelusuri aliran nilai pada pemanggilan fungsi dan menjelaskan perbedaan pass by value dan pass by reference serta jangkauan variabel lokal\\\garis
14 & Menulis fungsi rekursif dengan kasus basis dan langkah rekursif yang tepat, serta mendefinisikan \texttt{struct} dan memproses array of struct & Menelusuri pemanggilan rekursif dengan pohon pemanggilan atau tumpukan panggilan, dan menjelaskan peran kasus basis serta akses anggota struct\\\garis
15 & Menyelesaikan masalah yang menggabungkan array, string, fungsi, pengurutan atau pencarian, dan struct dalam satu program C++ dalam batas waktu ujian & Menelusuri program yang memakai array, fungsi, dan rekursi, lalu menjelaskan setiap perubahan nilai dengan istilah yang tepat\\\garis
16 & Menyelesaikan masalah komputasi dengan program C++ yang memakai array, string, fungsi, pengurutan, pencarian, rekursi, dan struct & Menelusuri dan menjelaskan perilaku program yang memakai array, fungsi, dan rekursi\\
\garisdasar
\end{longtable}}\par\vspace{6pt}
\Bagian{RPS}{Rencana pembelajaran mingguan}
Pertemuan reguler: review 15 menit, struggle 25 menit, materi 40 menit dengan pola tebak lalu buktikan, hands-on bertingkat 45 menit, Exit Ticket 25 menit. Bahan setiap pertemuan ada di folder \texttt{01 Slide}, \texttt{02 Hands-on}, dan \texttt{03 Modul}.
\par\vspace{4pt}\noindent{\fontsize{8.5pt}{11.5pt}\selectfont
\begin{longtable}{@{}>{\raggedright\arraybackslash}p{0.050\textwidth}>{\raggedright\arraybackslash}p{0.250\textwidth}>{\raggedright\arraybackslash}p{0.270\textwidth}>{\raggedright\arraybackslash}p{0.310\textwidth}>{\raggedright\arraybackslash}p{0.100\textwidth}@{}}
\kepala{P} & \kepala{Topik} & \kepala{Metode dan pasangan} & \kepala{Penilaian} & \kepala{Bobot}\\
\endhead
1 & Hello World dan Struktur Program C++ & Siklus 5 fase; Algoritma: Computational Thinking & Exit Ticket 1; Tugas 1: Kartu Perkenalan Digital & 3.87\%\\\garis
2 & Variabel, Tipe Data, Ekspresi, dan Input & Siklus 5 fase; Algoritma: Tipe Data, Variabel dan Ekspresi & Exit Ticket 2; Tugas 2: Program Pendaftaran Mahasiswa Baru & 3.87\%\\\garis
3 & Operator dan Precedence & Siklus 5 fase; Algoritma: Operator dan Precedence & Exit Ticket 3; Tugas 3: Kalkulator Lengkap dengan Validasi & 3.87\%\\\garis
4 & Percabangan & Siklus 5 fase; Algoritma: Percabangan & Exit Ticket 4; Tugas 4: Sistem Tiket Bioskop & 3.87\%\\\garis
5 & Perulangan Bagian 1 & Siklus 5 fase; Algoritma: Perulangan Bagian 1 & Exit Ticket 5; Tugas 5: Program Statistik Nilai & 3.87\%\\\garis
6 & Perulangan Bagian 2 dan Pola & Siklus 5 fase; Algoritma: Perulangan Bagian 2 dan Pattern & Exit Ticket 6; Tugas 6: Pattern Generator & 3.87\%\\\garis
7 & Review dan Simulasi UTS & Siklus 5 fase; Algoritma: Review dan Simulasi UTS & Exit Ticket 7 & 1.79\%\\\garis
8 & \textbf{UTS} (Pertemuan 1 sampai 6) & Ujian tertulis dan praktik, 120 menit & UTS: 2 soal A, 2 soal B & 25.00\%\\\garis
9 & Array Satu Dimensi & Siklus 5 fase; Algoritma: Array 1D & Exit Ticket 9; Tugas 9: Sistem Nilai Mahasiswa & 3.87\%\\\garis
10 & Array Dua Dimensi dan String & Siklus 5 fase; Algoritma: Array 2D dan String & Exit Ticket 10; Tugas 10: Sistem Nilai Kelas & 3.87\%\\\garis
11 & Searching & Siklus 5 fase; Algoritma: Searching & Exit Ticket 11; Tugas 11: Direktori Kontak & 3.87\%\\\garis
12 & Sorting & Siklus 5 fase; Algoritma: Sorting & Exit Ticket 12; Tugas 12: Sistem Leaderboard & 3.87\%\\\garis
13 & Fungsi dan Modularitas & Siklus 5 fase; Algoritma: Fungsi dan Modularitas & Exit Ticket 13; Tugas 13: Library Matematika & 3.87\%\\\garis
14 & Rekursi dan Struct & Siklus 5 fase; Algoritma: Rekursi dan Struct & Exit Ticket 14; Tugas 14: Perpustakaan Mini dengan Struct dan Rekursi & 3.87\%\\\garis
15 & Review dan Simulasi UAS & Siklus 5 fase; Algoritma: Review dan Simulasi UAS & Exit Ticket 15 & 1.79\%\\\garis
16 & \textbf{UAS} (Pertemuan 9 sampai 14) & Ujian tertulis dan praktik, 150 menit & UAS: 2 soal A, 2 soal B & 25.00\%\\
\garisdasar
\end{longtable}}\par\vspace{6pt}
Bobot per pertemuan dihitung dari 25\% Exit Ticket dibagi 14, ditambah 25\% tugas dibagi 12 bila ada tugas; UTS dan UAS masing-masing 25\%. Jumlah seluruh bobot tepat 100\%.

\Bagian{RPS}{Penilaian}
\par\vspace{4pt}\noindent{\fontsize{9.5pt}{13pt}\selectfont
\begin{tabular}{@{}>{\raggedright\arraybackslash}p{0.150\textwidth}>{\raggedright\arraybackslash}p{0.380\textwidth}>{\raggedright\arraybackslash}p{0.120\textwidth}>{\raggedright\arraybackslash}p{0.160\textwidth}>{\raggedright\arraybackslash}p{0.160\textwidth}@{}}
\kepala{Komponen} & \kepala{Pelaksanaan} & \kepala{Bobot} & \kepala{CPMK0607} & \kepala{CPMK0608}\\
Exit Ticket & 14 kali (P1–P7, P9–P15), rata-rata & 25\% & 12,5 & 12,5\\\garis
Tugas & 12 kali (P1–P6, P9–P14), rata-rata & 25\% & 12,5 & 12,5\\\garis
UTS & P8, 4 soal, 120 menit & 25\% & 12,5 & 12,5\\\garis
UAS & P16, 4 soal, 150 menit & 25\% & 12,5 & 12,5\\\garis
\textbf{Total} &  & \textbf{100\%} & \textbf{50} & \textbf{50}\\
\garisdasar
\end{tabular}}\par\vspace{6pt}
Setiap instrumen dibagi rata menjadi Bagian A (implementasi kode, 50\% poin, CPMK0607) dan Bagian B (tracing dan penjelasan, 50\% poin, CPMK0608). Karena itu kontribusi setiap CPMK pasti 50 poin, apa pun bobot antarkomponennya.

Nilai minimum lulus setiap CPMK adalah 50 dari skala 0–100 pada CPMK itu. Mahasiswa dengan salah satu CPMK di bawah 50 wajib mengikuti remedial.

\Sub{Rubrik umum empat level}
\par\vspace{4pt}\noindent{\fontsize{9.5pt}{13pt}\selectfont
\begin{tabular}{@{}>{\raggedright\arraybackslash}p{0.160\textwidth}>{\raggedright\arraybackslash}p{0.140\textwidth}>{\raggedright\arraybackslash}p{0.680\textwidth}@{}}
\kepala{Level} & \kepala{Skor} & \kepala{Deskripsi}\\
Sangat baik & 85–100 & Memenuhi semua kriteria dengan tepat dan alasan yang jelas\\\garis
Baik & 70–84 & Memenuhi kriteria dengan satu kekurangan kecil\\\garis
Cukup & 55–69 & Memenuhi sebagian kriteria; ada kesalahan yang memengaruhi hasil\\\garis
Kurang & di bawah 55 & Sebagian besar kriteria tidak terpenuhi, atau hasil disalin\\
\garisdasar
\end{tabular}}\par\vspace{6pt}
Setiap tugas memakai lima kriteria: A1 program berjalan (20), A2 dan A3 kriteria isi (15 dan 15), B1 penjelasan (30), dan B2 cerita error dan pengujian (20). Nilai kriteria = poin × skor level. Rincian per tugas ada di modul pertemuannya.

\Sub{Konversi nilai akhir}
\par\vspace{4pt}\noindent{\fontsize{9.5pt}{13pt}\selectfont
\begin{tabular}{@{}>{\raggedright\arraybackslash}p{0.200\textwidth}>{\raggedright\arraybackslash}p{0.300\textwidth}@{}}
\kepala{Huruf} & \kepala{Rentang}\\
A & ≥ 75\\\garis
AB & 70–74,99\\\garis
B & 65–69,99\\\garis
BC & 60–64,99\\\garis
C & 50–59,99\\\garis
D & 40–49,99\\\garis
E & < 40\\
\garisdasar
\end{tabular}}\par\vspace{6pt}
Sesuai A-03 Standar Penilaian Pembelajaran ITERA.

\Bagian{RPS}{Referensi}
\begin{enumerate}
\item Deitel, P. dan Deitel, H. C++ How to Program, edisi ke-10. Pearson.
\item Stroustrup, B. Programming: Principles and Practice Using C++, edisi ke-2. Addison-Wesley.
\item Gaddis, T. Starting Out with C++. Pearson.
\item learncpp.com dan cppreference.com.
\end{enumerate}
\end{document}
